\section{Linear Algebra Review}

	In functional analysis, there are two big theorems that we care about, and those are duality theorems. Let's review some linear algebra. First, the \textit{Riesz Representation Theorem}. Let $X$ be a Hilbert space over $\mathbb{C}$. Then, for any linear functional $l \in X^*$, there exists a unique vector $v \in X$ such that 
	\[l(x) = \langle x, v \rangle_X\]
	for all $x \in X$. That is, any linear functional $l(\cdot)$ can be represented as an inner product of the \textit{first argument} $\langle \cdot , v \rangle$. Remember that over $\mathbb{C}$, the inner product has first argument linearity and second argument conjugate linearity. This map $l \mapsto v$ is a \textit{natural isomorphism}, but only under the condition that $X$ is a Hilbert space. This does not work when $X$ does not have an inner product. 

	Let $\mathcal{L}(X, Y)$ be the vector space of linear maps from vector space $X$ to vector space $Y$. Then, the adjoint map 
	\[\ast : \mathcal{L}(X, Y) \longrightarrow \mathcal{L}(Y^*, X^*)\]
	is an isometric isomorphism. We construct this isomorphism as such: Given a fixed linear map $T \in \mathcal{L}(X, Y)$, choose some vector $x \in X$. The dual spaces of $X, Y$ automatically exist without invoking anything else. Now choose any vector $l \in Y^*$. Then, there exists a unique $T^* \in \mathcal{L}(Y^*, X^*)$ such that 
	\[l (Tx) = (T^* l) (x)\]
	for all $x \in X, l \in X^*$. We can visualize it with the following diagram. 
	\[\begin{tikzcd}
	x \in X \arrow{r}{T} & Y \ni Tx \\
	T^* l \in X^* & \arrow{l}{T^*} Y^* \ni l
	\end{tikzcd}\]
	This is called the \textit{Banach adjoint}. If we assume even further that $X$ and $Y$ are Hilbert spaces, then we can invoke Riesz representation theorem to define an isomorphism $i_Y: Y \rightarrow Y*$ and $i_X: X \rightarrow X^*$ such that for any given $x \in X$ and $y \in Y$, 
	\begin{align*}
			i_X (x) = l_x & \text{ s.t. } l_x (\cdot) = \langle \cdot, x \rangle_X \\ 
			i_Y (y) = l_y & \text{ s.t. } l_y (\cdot) = \langle \cdot y \rangle_Y 
	\end{align*}
	and define the \textit{Hilbert adjoint} to be, $T^H = i_X^{-1} \circ T^* \circ i_Y: Y \longrightarrow X$, which satisfies 
	\[\langle y, Tx \rangle_Y = \langle T^* y, x \rangle_X\]
	We differentiate between these two definitions because some texts define the adjoint to be the map from $Y$ to $X$ rather than their duals, but with the Riesz isometry, the Banach and Hilbert adjoints are precisely the same. The big idea is that $\ast$ is an isometric isomorphism. We will prove only linearity. Assume $A, B \in \mathcal{L}(X, Y)$. Then, we wish to show that $\ast(A + B) = \ast(A) + \ast(B)$. Given $A + B \in \mathcal{L}(X, Y)$, $\ast(A + B)$ is defined such that for all $x \in X, l \in Y^*$, 
	\begin{align*}
			[ (A + B)^* (l) ] (x) & = (l) [(A + B)(x)] \\
			& = (l) [A x + Bx] & (\text{by definition of } A + B) \\ 
			& = l(Ax) + l(Bx) & (\text{by linearity of } l) \\
			& = (A^* l)(x) + (B^* l)(x) & (\text{definition of adjoint}) 
	\end{align*}
	We can also see that the map $T^\prime = i_Y \circ T \circ i_X^{-1} : X^* \longrightarrow Y^*$ is also well-defined. This is not necessarily the inverse of $T^*$. 

	So if we have a Hilbert space $\mathcal{H}$, the adjoint operator is an isometric isomorphism from $\mathcal{L}(\mathcal{H}, \mathcal{H})$ to itself (since by Riesz rep theorem $\mathcal{H}^* \simeq \mathcal{H}$). 
	\[\ast : \mathcal{L}(\mathcal{H}, \mathcal{H}) \longrightarrow \mathcal{L}(\mathcal{H}, \mathcal{H})\]
	So therefore $\mathcal{L}(\mathcal{H}, \mathcal{H})$ has the following properties. Let $A: \mathcal{H} \longrightarrow \mathcal{H}$ and $B: \mathcal{H} \longrightarrow \mathcal{H}$ be elements of $\mathcal{L}(\mathcal{H}, \mathcal{H})$. 
	\begin{enumerate}
			\item As we know, $\mathcal{L}(\mathcal{H}, \mathcal{H})$ is a vector space since $A + B, \lambda A \in \mathcal{L}(\mathcal{H}, \mathcal{H})$ for any $\lambda \in \mathbb{C}$. 
			\item We can define multiplication on $\mathcal{L}(\mathcal{H}, \mathcal{H})$ to be simply the composition of operators, so $AB \in \mathcal{L}(\mathcal{H}, \mathcal{H})$. 
			\item We can define the adjoint operator on $\mathcal{L}(\mathcal{H}, \mathcal{H})$ from above, which satisfies $A^{**} = A$ and $(AB)^* = B^* A^*$. 
	\end{enumerate}
	The first two conditions simply make $\mathcal{L}(\mathcal{H}, \mathcal{H})$ an algebra (in the algebraic sense), but the third condition makes it a $C^*$-algebra. 

	Now given the $L^2$ space (of complex-valued functions over a measure space $X$), we can upgrade it from Banach to Hilbert by defining the inner product 
	\begin{equation}
		\langle f, g \rangle \coloneqq \int_X \overline{f(x)} g(x) \, d\mu
	\end{equation}
	Now, given $p, q$ satisfying $\frac{1}{p} + \frac{1}{q} = 1$, let $f \in L^p$ and $g \in L^q$. Then, $f$ defines a linear map 
	\begin{equation}
		T_f: g \mapsto \int_X f g \, d\mu 
	\end{equation}
	Sometimes, this map $T_f$ may be written $\int_X f$, but this does not denote an integral; it is purely notation. $\int_X f$ does not even have to be finite since $f$ is not $L^1$. They are only integrals when you feed them the appropriate terms. 

	Some definitions. 
	\begin{enumerate}
		\item $c_0 = \big\{ \{a_j\} \,|\, a_j \rightarrow 0 \big\}$ is the set of convergent sequences. It can be seen that $c_0 \subset \ell_\infty$ since every convergent sequence is bounded. 
		\item $\ell_1 = \big\{ \{a_j\} \,|\, \sum_j |a_j| < \infty\}$ is the set of sequences such that the sum of all (the absolute value of) its terms is finite. But upon looking at this, we can see that $\ell_1$ is precisely the $L^1$ space of functions over the measure space $(\mathbb{N}, 2^\mathbb{N}, \mu)$, where $\mu$ is the counting measure (where the measure of a subset is determined by the number of elements in that set). Note that $\ell_1 \subset c_0$. 
		\item $\ell_2 = \big\{ \{a_j\} \,|\, \sum a_j^2 < \infty \big\}$. That is, this is the $L^2 (\mathbb{N})$ space under the counting measure. 
		\item $\ell_p$ can be defined as the same space under the counting measure, but equipped with the $p$-norm. 
		\item $\ell_\infty = \big\{ \{ a_j \} \,|\, |a_j| \leq c \text{ for all} j \big\}$ for some constant $c$. Clearly, this is the $L^\infty$-norm, and it must be for all $j$ since the counting measure on any nonempty subset of $\mathbb{N}$ is at least $1$.  
	\end{enumerate}
	In summary, 
	\[\ell_1 \subset \ell_2 \subset \ell_3 \subset \ldots \subset c_0 \subset \ell_\infty\]
	To prove this, let's assume that some sequence $f \in \ell_p$, and let $1 \leq p < q \leq \infty$. Then,$f \in c_0$ and therefore is convergent to $0$. There must be some $K \in \mathbb{N}$ such that $f(k) < 1$ for $k > K$. So, for all $k > K$, we have $|f(k)|^q < |f(k)|^p < 1$, and so 
	\[\sum_{k > K} |f(k)|^q < \sum_{k > K} |f(k)|^p\]
	The first $k$ terms would not affect finiteness of the norm so 
	\[||f||_p = \sum_{k \in \mathbb{N}} |f(k)|^p < \infty \implies ||f||_q = \sum_{k \in \mathbb{N}} |f(k)|^q < \infty\]
	and so $\ell_p \subset \ell_q$. \qed

	Let's talk about duality. Unlike that of a finite-dimensional vector space, the dual of an infinite dimensional space is not isomorphic to itself. This dual space just some abstract space of linear functionals, but we can informally think of it as sort of an "infinite dot product." That is, given a sequence $\{ f(n) \}_{n \in \mathbb{N}}$, an element $\phi$ in its dual space would act on $f$ as 
	\[\phi(f) \coloneqq \sum_{n \in \mathbb{N}} f(n) \phi(n) \in \mathbb{C}\]
	For example, the dual of $c_0$ is not itself because $f = \big(\frac{1}{\sqrt{n}}\big)_n \in c_0$, but the functional $\phi = \big(\frac{1}{\sqrt{n}} \big)$ acting on $f$ would be infinite. 
	\[\phi(f) = \sum_{n \in \mathbb{N}} \frac{1}{\sqrt{n}} \frac{1}{\sqrt{n}} = \sum_{n \in \mathbb{N}} \frac{1}{n} = \infty \not\in \mathbb{C}\]
	So what is the dual space of $c_0$? It turns out that $c_0^* \simeq \ell_1$. In order to prove this, we must prove that the map $\phi: \ell_1 \longrightarrow c_0^*$ defined 
	\[\phi(g) (f) \coloneqq \sum_{n \in \mathbb{N}} g(n) \, f(n)\]
	is an isomorphism. Linearity is easy since 
	\begin{align*}
			\phi(\lambda g_1 + g_2) (f) & = \sum_{n \in \mathbb{N}} (\lambda g_1 + g_2) (n) \, f(n) \\
			& = \sum \big( \lambda g_1 (n) + g_2 (n) \big) \, f(n) \\
			& = \sum \lambda g_1(n) \, f(n) + g_2 (n) \, f(n) \\
			& = \lambda \sum g_1(n) \, f(n) + \sum g_2(n) \, f(n) \\
			& = \lambda \phi(g_1) (f) + \phi (g_2) (f)
	\end{align*}
	Now, we want to prove boundedness (i.e. $\phi(g)$ is bounded for all $g$ so doesn't explode to $\infty$). By Holder's Inequality, we have $f \in c_0 \subset \ell_\infty$ and $g \in \ell_1$, so 
	\[\sum_{n \in \mathbb{N}} |f(n) \, g(n)| \leq ||f||_\infty ||g||_1 < \infty\]
	since $\phi(g)(f)$ converges absolutely, it also converges. To prove injectivity, we can just prove that the nullspace of $\phi$ is trivial. Assume that a nontrivial $l \in \ell_1$ is in the nullspace. Then, there exists a $k$th nonzero term: $l(k) \neq 0$. But we have claimed that $\phi(l)(f) = 0$ for all $f \in c_0$, but setting $f(k) = 1$, this makes a contradiction. To prove surjectivity, let $\varphi \in c_0^*$ be a linear functional. Then, we can construct a sequence $g$ element-wise by assigning 
	\[g_n = \varphi(e_n)\]
	where $e_n$ is the $n$th basis element of $c_0$. Then, $g = (g_n)_{n \in \mathbb{N}}$ is a sequence, and we want to show that it is in $\ell_1$. 

	Furthermore, given that $p, q$ are duals, i.e. $\frac{1}{p} + \frac{1}{q} = 1$, we have 
	\begin{enumerate}
			\item $c_0^* = \ell_1$
			\item $\ell_1^* = \ell_\infty$ 
			\item $\ell_2^* = \ell_2$. Since the dual space of $\ell_2$ is itself, using Riesz Representation theorem we can construct an inner product on $\ell_2$ to make it a Hilbert space. $\ell_2$ is the only space that can be construct into a Hilbert space. The rest are Banach spaces. 
	\end{enumerate}

